\documentclass[12pt]{article}
\usepackage[utf8]{inputenc}
\usepackage[T1]{fontenc}
\usepackage{hyperref}
\title{Tektiteko Language Resources for K-12 Curriculum Development}
\author{Compiled by Claude Sonnet 5 (claude-sonnet-5) for Andr\'as Kornai}
\date{}
\begin{document}
\maketitle

\section{Introduction}
Tektiteko (also Tektitek, Tectiteco or Teko; ISO 639-3 ttc) is a Mayan language of the Mamean branch. It is spoken in southern Huehuetenango, Guatemala, and in Chiapas, Mexico, and is very closely related to Mam. Speaker counts in the search results vary from about 3,000 to 5,900 depending on source and date. Little material was found online.

\section{Primary Educational Resources}
No primers or textbooks were found in search results. Omniglot states that books exist, including grammars, dictionaries and textbooks. I could not confirm this independently, and no specific titles were found.

\section{Linguistic and Reference Materials}
Wikipedia cites a grammar and idioms work by Paul Stevenson (1986--1987) and Terrence Kaufman's 1969 \emph{Teco -- A New Mayan Language}; I did not see these items themselves.
\begin{itemize}
\item Wikipedia: \url{https://en.wikipedia.org/wiki/Tektitek_language}
\item Omniglot: \url{https://mail.omniglot.com/writing/tektitek.htm}
\end{itemize}

\section{Religious and Community Resources}
None found.

\section{Digital and Audio Resources}
\begin{itemize}
\item Native Languages: vocabulary, pronunciation guide and a Spanish--Tektiteko glossary with some audio. \url{https://www.native-languages.org/tectiteco.htm}
\item Endangered Languages Project entry, with a bibliography. \url{https://endangeredlanguages.com/node/89067}
\item AILLA (UT Austin) Mayan collection page, which lists Tektiteko. \url{https://sites.utexas.edu/ailla/mayan/}
\end{itemize}

\section{Curriculum Development Context}
Omniglot says the language is taught in primary schools. The reference description says it is not known to be taught, and I found nothing to settle the conflict. Because Tektiteko is close to Mam, Mam-language educational materials may be adaptable, but I did not research them.
\end{document}